% =========================================================
%  sesi/03/sesi.tex
%  Sesi 3 — Pemain bergerak dengan vektor dan delta time
% =========================================================

\sesi{03}{Pemain Bergerak dengan Vektor dan Delta Time}{1.2}{120 menit}

Pada sesi 2 siklus permainan sudah berputar, tetapi belum ada yang berpindah
tempat. Layar hanya menampilkan latar yang sama dari bingkai ke bingkai. Pada
minggu 3 ini tahap pembaruan akhirnya diberi pekerjaan: seorang pemain
berukuran 36 piksel muncul di tengah jendela dan bergerak ke kiri atau ke kanan
selama tombol arah ditahan.

Gerak itu tidak dihitung dari jumlah bingkai, melainkan dari waktu. Posisi
pemain disimpan sebagai vektor (\istilah{Vector2}) bertipe pecahan, dan setiap
bingkai posisinya ditambah kecepatan dikali selang waktu (\istilah{delta
time}). Dengan cara itu, pemain menempuh jarak yang sama dalam satu detik, baik
program berjalan pada 30 maupun 60 \istilah{FPS}. Inilah pembeda penting antara
program yang hanya berkedip dan permainan yang geraknya dapat dipercaya.

% ---------------------------------------------------------
\bagian{Tujuan Sesi}

\begin{tujuanbox}
Setelah menuntaskan sesi ini, Anda dapat menunjukkan:
\begin{enumerate}[leftmargin=1.5em, itemsep=0.2em]
  \item pemain berukuran 36 $\times$ 36 piksel yang tampil di tengah jendela
        \NamaPermainan;
  \item pemain yang bergerak ke kiri dan ke kanan selama tombol panah atau
        tombol A dan D ditahan;
  \item kecepatan 260 piksel per detik yang tetap sama ketika program
        dijalankan pada 30 FPS dan 60 FPS;
  \item posisi pemain yang tersimpan sebagai pecahan, sementara
        \istilah{rect} untuk menggambar diselaraskan dengan pembulatan;
  \item perhitungan $\vec{p}_{\text{baru}} = \vec{p} + \vec{v}\,\Delta t$
        beserta arti setiap lambangnya.
\end{enumerate}
\end{tujuanbox}

% ---------------------------------------------------------
\bagian{Prasyarat}

\begin{prasyaratbox}
\begin{itemize}[leftmargin=1.4em, itemsep=0.25em]
  \item Checkpoint sesi 2 pada \file{code/sesi-02} sudah berjalan: siklus
        permainan berputar, FPS tampil pada judul, dan jendela dapat ditutup
        dengan rapi.
  \item Anda telah membaca Bab V buku ajar, khususnya bagian posisi sebagai
        vektor, pembacaan keadaan tombol, gerak berbasis selang waktu, dan
        penyelarasan posisi dengan \code{Rect}.
  \item Anda masih mengingat perbedaan peristiwa dan keadaan tombol dari
        Bab IV; perbedaan itu menjadi dasar pemilihan
        \code{pygame.key.get_pressed()} pada sesi ini.
  \item Anda dapat menjalankan program dengan pilihan baris perintah, misalnya
        \code{python main.py --fps 30}.
\end{itemize}
\end{prasyaratbox}

% ---------------------------------------------------------
\bagian{Peta Kemampuan}

\petasesi
  {Siklus berjalan dan FPS tampil pada judul, tetapi layar hanya berisi latar
   polos.}
  {Pemain 36 piksel, posisi \code{Vector2}, \code{key.get_pressed()},
   kecepatan 260 piksel per detik, dan penyelarasan \code{rect}.}
  {Percepatan, gravitasi, lompatan, dan pantulan yang menambah kecepatan
   sebelum posisi diperbarui pada sesi 4.}

% ---------------------------------------------------------
\bagian{Konsep yang Dipakai}

\subsection*{Vektor menyatukan dua komponen}

Posisi pada bidang datar memiliki dua angka: jarak mendatar dan jarak tegak.
Menyimpannya pada dua variabel terpisah membuat setiap operasi harus ditulis
dua kali dan membuat kekeliruan menukar komponen sulit terlihat.
\code{pygame.Vector2} menyatukan keduanya menjadi satu nilai, sehingga
penjumlahan dan perkalian dengan satu angka bekerja per komponen secara
otomatis. Perlu diingat bahwa pada koordinat layar sumbu tegak menuju ke bawah;
gerak ke atas berarti mengurangi nilai $y$, bukan menambahnya.

\subsection*{Keadaan tombol bukan peristiwa}

Peristiwa (\istilah{event}) \code{KEYDOWN} hanya dilaporkan sekali pada saat
tombol mulai ditekan. Peristiwa itu tidak cocok untuk gerak yang berlangsung
terus, karena pemain yang menahan tombol hanya akan melihat satu perpindahan
kecil. Untuk keperluan itu Pygame menyediakan
\code{pygame.key.get_pressed()}, yang mengembalikan keadaan seluruh tombol pada
bingkai itu. Selama tombol masih ditahan, nilainya benar; begitu dilepas,
nilainya salah. Inilah cara membaca masukan yang dipakai sesi ini.

\subsection*{Perpindahan dari kecepatan dan selang waktu}

Kecepatan dinyatakan dalam piksel per detik, sedangkan selang waktu dinyatakan
dalam detik. Perpindahan satu bingkai adalah hasil kalinya:

\[
\vec{p}_{\text{baru}} = \vec{p} + \vec{v}\,\Delta t.
\]

Lambang $\vec{p}$ menyatakan posisi sebelum pembaruan, $\vec{v}$ menyatakan
kecepatan dalam piksel per detik, dan $\Delta t$ menyatakan selang waktu satu
bingkai dalam detik. Hasil perkalian $\vec{v}\,\Delta t$ bersatuan piksel,
sehingga dapat langsung dijumlahkan dengan posisi. Nilai $\Delta t$ diperoleh
dari pengendali waktu:

\[
\Delta t = \frac{\text{tick}}{1000}.
\]

Lambang \text{tick} menyatakan nilai yang dikembalikan
\code{clock.tick(FPS)} dalam milidetik, dan angka 1000 mengubahnya menjadi
detik. Apabila selang waktu dua kali lebih panjang, perpindahannya dua kali
lebih jauh, tetapi jarak yang ditempuh dalam satu detik tetap sama.

\subsection*{Posisi pecahan, rect bulat}

Posisi yang sebenarnya harus disimpan sebagai pecahan agar perpindahan kecil
tetap terakumulasi. Pada kecepatan 260 piksel per detik dan 60 FPS,
perpindahan setiap bingkainya sekitar 4,33 piksel; pada kecepatan rendah,
pecahan itu jauh lebih penting. \code{Rect} tidak menyimpan pecahan karena
letak piksel pada layar memang bilangan bulat. Karena itu, setiap bingkai
\code{rect} diselaraskan dari posisi pecahan dengan \code{round()}, bukan
dipotong dengan \code{int()}. Membulatkan memilih bilangan bulat terdekat,
sedangkan memotong selalu bergerak ke arah nol dan membuat pemain tampak
tertinggal hampir satu piksel.

% ---------------------------------------------------------
\bagian{Tutorial}

Kerjakan langkah berikut berurutan. Setiap langkah selesai ketika hasilnya
dapat Anda lihat, bukan ketika kodenya baru diketik.

\begin{langkah}
  \item \textbf{Siapkan folder sesi.} Salin folder \file{code/sesi-02} menjadi
        folder kerja baru, lalu buka di editor dan terminal. Anda berangkat
        dari siklus yang sudah berjalan dan berjudul FPS.

  \item \textbf{Tambahkan pengaturan pemain.} Pada
        \file{gudang_artefak/settings.py}, tambahkan \code{PLAYER_SIZE = 36},
        \code{PLAYER_SPEED = 260.0}, dan \code{PLAYER_COLOR}. Beri keterangan
        bahwa kecepatan memakai satuan piksel per detik, bukan piksel per
        bingkai.

  \item \textbf{Buat berkas benda.} Buat \file{gudang_artefak/entities.py}
        dan impor \code{pygame} serta \code{settings}. Berkas ini akan memuat
        kelas \code{Player}; memisahkannya dari \code{game.py} membuat aturan
        gerak mudah dicari ketika program bertambah besar.

  \item \textbf{Tulis kelas pemain.} Buat \code{Player} sebagai turunan
        \code{pygame.sprite.Sprite}. Di \code{__init__}, buat permukaan
        berukuran 36 $\times$ 36, isi dengan warna pemain, lalu buat
        \code{self.rect} dari permukaan itu. Simpan posisi sebenarnya pada
        \code{self.position = pygame.Vector2(x, y)}.

  \item \textbf{Tulis pembaruan pemain.} Buat metode \code{update(delta_seconds,
        direction)}. Tambahkan \code{direction * S.PLAYER_SPEED *
        delta_seconds} ke \code{self.position.x}, jepit nilainya agar pemain
        tidak keluar jendela, lalu selaraskan \code{self.rect.x} dengan
        \code{round(self.position.x)}.

  \item \textbf{Tempatkan pemain di tengah.} Di \code{Game.__init__}, buat
        \code{Player} pada posisi setengah lebar dan setengah tinggi jendela
        dikurangi setengah ukuran pemain. Sesuaikan \code{player.bounds}
        dengan ukuran jendela yang benar-benar dipakai, bukan hanya dengan
        nilai bawaan.

  \item \textbf{Hitung selang waktu.} Di awal \code{run()}, ubah
        \code{self.clock.tick(self.args.fps)} menjadi
        \code{delta_seconds = self.clock.tick(self.args.fps) / 1000.0}.
        Pembagian seribu mengubah milidetik menjadi detik.

  \item \textbf{Panggil pembaruan.} Di dalam \code{run()}, susun urutannya
        menjadi \code{handle_events()}, \code{update(delta_seconds)},
        \code{draw()}, lalu \code{pygame.display.flip()}.

  \item \textbf{Baca keadaan tombol.} Buat \code{Game.update(delta_seconds)}.
        Di dalamnya panggil \code{keys = pygame.key.get_pressed()}, lalu
        hitung \code{direction = int(keys[pygame.K_RIGHT] or
        keys[pygame.K_d]) - int(keys[pygame.K_LEFT] or keys[pygame.K_a])}.
        Teruskan \code{direction} ke \code{self.player.update()}.

  \item \textbf{Gambar pemain.} Di \code{draw()}, setelah mengisi latar,
        tempelkan \code{self.player.image} pada \code{self.player.rect}
        dengan \code{self.screen.blit()}.
\end{langkah}

\begin{ujibox}
Setelah Langkah 10, periksa empat hal berikut.
\begin{itemize}[leftmargin=1.4em, itemsep=0.2em]
  \item Pemain tampil sebagai kotak 36 $\times$ 36 piksel di tengah jendela,
        bukan di pojok kiri atas.
  \item Menahan panah kanan atau D menggerakkan pemain ke kanan selama tombol
        ditahan, bukan satu langkah kecil lalu berhenti.
  \item Menahan panah kiri atau A menggerakkan pemain ke kiri, dan pemain
        berhenti di tepi jendela tanpa keluar dari layar.
  \item Pemain tidak meninggalkan jejak. Bila ada jejak, periksa urutan
        \code{fill}, \code{blit}, lalu \code{flip}.
\end{itemize}
\end{ujibox}

Sekarang uji ketahanan gerak terhadap perubahan laju bingkai. Jalankan
\code{python main.py --fps 60}, lalu \code{python main.py --fps 30}. Ukur waktu
yang diperlukan pemain untuk melintasi hampir seluruh lebar jendela, misalnya
dari tepi kiri ke tepi kanan, dengan pengukur waktu di ponsel. Pada kedua
pengaturan, waktunya harus kira-kira sama, yaitu sekitar 3,5 detik untuk jarak
900 piksel pada kecepatan 260 piksel per detik.

\begin{ujibox}
Tunjukkan tiga bukti berikut sebelum melanjutkan.
\begin{itemize}[leftmargin=1.4em, itemsep=0.2em]
  \item Waktu lintasan pada 30 FPS dan 60 FPS tidak berbeda jauh; yang berubah
        hanyalah jumlah bingkai, bukan jarak per detik.
  \item Gerak pada 30 FPS terlihat lebih tersendat, tetapi pemain tetap
        menempuh jarak yang sama dalam satu detik.
  \item Bila \code{/ 1000.0} dihapus sementara, pemain bergerak sangat cepat
        dan menempel di tepi jendela. Kembalikan pembagian itu setelah
        percobaan.
\end{itemize}
\end{ujibox}

% ---------------------------------------------------------
\bagian{Kode Sesi Ini}

Berkas berikut adalah isi checkpoint sesi 3. Bagian yang berubah dari sesi 2
adalah \file{settings.py} (pengaturan pemain), \file{entities.py} (kelas
\code{Player}), \file{game.py} (pembaruan posisi dan penggambaran pemain), dan
\file{main.py} (tetap sama, tetapi pilihan \code{--fps} kini mengubah gerak).
Ketiklah sendiri, lalu bandingkan hasilnya.

\tampilkodebaris[Python]{code/sesi-03/gudang_artefak/settings.py}{10}{22}
  {settings.py: ukuran, warna, dan kecepatan pemain}

\tampilkodebaris[Python]{code/sesi-03/gudang_artefak/entities.py}{14}{39}
  {entities.py: posisi pecahan, pembaruan, dan penyelarasan rect}

\tampilkodebaris[Python]{code/sesi-03/gudang_artefak/game.py}{19}{31}
  {game.py: pemain dibuat sekali di tengah jendela}

\tampilkodebaris[Python]{code/sesi-03/gudang_artefak/game.py}{33}{56}
  {game.py: selang waktu dihitung dari clock.tick dan diteruskan ke pembaruan}

\tampilkodebaris[Python]{code/sesi-03/gudang_artefak/game.py}{67}{77}
  {game.py: pembacaan keadaan tombol dan penggambaran pemain}

\checkpoint{03}{Pemain 36 piksel bergerak mendatar dengan kecepatan 260 piksel
per detik, kecepatannya sama pada 30 dan 60 FPS, dan posisi pecahannya
diselaraskan dengan \code{rect} melalui pembulatan.}

% ---------------------------------------------------------
\bagian{Kesalahan Umum}

\begin{galatbox}
\begin{itemize}[leftmargin=1.4em, itemsep=0.35em]
  \item \textbf{Pemain tidak bergerak sama sekali.}
        Gejala: pemain tampil, tetapi tombol tidak berpengaruh. Periksa apakah
        \code{Game.update()} benar-benar dipanggil dari \code{run()} dan apakah
        \code{delta_seconds} diteruskan ke \code{player.update()}.
  \item \textbf{Pemain hanya bergerak satu langkah setiap penekanan.}
        Gejala: pemain melompat sedikit lalu diam walaupun tombol ditahan.
        Periksa pembacaan masukan; \code{KEYDOWN} hanya terjadi sekali,
        sedangkan \code{pygame.key.get_pressed()} membaca keadaan tombol.
  \item \textbf{Pemain bergerak terlalu cepat.}
        Gejala: pemain melesat dan menempel di tepi jendela. Periksa pembagian
        \code{/ 1000.0}; memakai nilai milidetik sebagai detik membuat gerak
        seribu kali lebih cepat.
  \item \textbf{Pemain meninggalkan jendela.}
        Gejala: pemain menghilang di tepi. Periksa penjepitan posisi dan
        pastikan \code{player.bounds} memakai ukuran jendela yang sebenarnya,
        bukan hanya nilai bawaan.
  \item \textbf{\code{AttributeError: 'tuple' object has no attribute 'x'}.}
        Gejala: program berhenti pada baris \code{position.x}. Periksa apakah
        posisi dibuat dengan \code{pygame.Vector2}, bukan dengan pasangan
        angka biasa.
  \item \textbf{Gerak tersendat pada kecepatan rendah.}
        Gejala: pemain diam beberapa bingkai lalu melompat satu piksel.
        Periksa penyimpanan posisi; bila posisi disimpan sebagai bilangan
        bulat, pecahan kecilnya terpotong dan hilang.
\end{itemize}
\end{galatbox}

% ---------------------------------------------------------
\bagian{Lembar Kerja}

Catat hasil pengukuran Anda. Isi kolom terakhir setelah Anda membandingkan
kedua laju bingkai dan mencoba satu nilai kecepatan yang lebih rendah.

{\small
\renewcommand{\arraystretch}{1.4}
\begin{tabularx}{\linewidth}{@{}C{0.8cm} L{5.2cm} Y@{}}
\toprule
\textbf{No} & \textbf{Hasil yang dicatat} & \textbf{Keterangan Anda} \\
\midrule
1 & Kecepatan yang ditetapkan dan ukuran pemain & \\[0.55cm]
2 & Waktu lintasan pada 30 FPS dan 60 FPS & \\[0.9cm]
3 & Jarak yang ditempuh dalam satu detik pada kedua laju & \\[0.9cm]
4 & Contoh nilai posisi pecahan dan nilai \code{rect} yang tampil & \\[0.9cm]
5 & Perilaku pemain ketika kecepatan diturunkan menjadi 20 piksel per detik & \\[0.9cm]
6 & Satu kendala yang Anda hadapi dan cara menyelesaikannya & \\[0.9cm]
\bottomrule
\end{tabularx}
}

% ---------------------------------------------------------
\bagian{Checklist Sesi}

\begin{checklistbox}
\begin{checklist}
  \item Saya dapat menunjukkan pemain 36 piksel yang tampil di tengah jendela.
  \item Saya dapat menggerakkan pemain ke kiri dan ke kanan dengan tombol
        panah maupun tombol A dan D.
  \item Saya dapat menunjukkan bahwa kecepatan pemain tetap sama pada 30 FPS
        dan 60 FPS.
  \item Saya dapat menjelaskan arti setiap lambang pada
        $\vec{p}_{\text{baru}} = \vec{p} + \vec{v}\,\Delta t$ dan
        $\Delta t = \text{tick}/1000$.
  \item Saya dapat menunjukkan posisi pecahan pemain dan \code{rect} bulat
        yang diselaraskan dengan \code{round()}.
  \item Saya dapat menjelaskan mengapa posisi disimpan sebagai pecahan,
        sedangkan \code{rect} dipakai untuk menggambar.
\end{checklist}
\end{checklistbox}

% ---------------------------------------------------------
\bagian{Pengait ke Sesi Berikutnya}

Pemain pada sesi ini bergerak lurus dengan kecepatan tetap; begitu tombol
dilepas, kecepatannya langsung nol. Sebagian besar permainan menuntut lebih:
benda yang jatuh makin cepat, bola yang memantul kehilangan sebagian
kecepatannya, dan karakter yang melompat melambat di dekat titik tertinggi.
Semua itu berasal dari percepatan. Pada sesi 4 Anda menambahkan
$\vec{v}_{\text{baru}} = \vec{v} + \vec{a}\,\Delta t$ sebelum posisi
diperbarui. Siklus dan penyelarasan \code{rect} tidak berubah; yang bertambah
hanyalah satu tingkat perhitungan.
